\section{主理想整环中的不可约因子分解}

\begin{definition}{不可约元代表系}{}
	设$A$为整环,$\left(p_{i}\right)_{i\in I}$是$A$的不可约元族.若对$A$的每个不可约元$p$,存在唯一的$i\in I$使得$p$和$p_{i}$相伴,则称$\left(p_{i}\right)_{i\in I}$是不可约元代表系.
\end{definition}

\begin{lemma}{每个左理想都有限生成的环一定左Noether}{}
	设$A$是环.如果$A$的每个左理想都有限生成,则$A$中的左理想构成的每个非空集合$\Phi$在包含偏序下一定有极大元.
\end{lemma}

\begin{theorem}{PID环的唯一分解性}{}
	设$A$是PID,$\left(p_{i}\right)_{i\in I}$是$A$的不可约元代表系,则对每个$x\in A\setminus\left\{0\right\}$,存在唯一的$c\in A^{\times}$和$\left(n_{i}\right)_{i\in I}\in\N^{\left(I\right)}$使得$x=c\prod\limits_{i\in I}p_{i}^{n_{i}}$.
\end{theorem}

\begin{remark}{}{}
	固定$x\in A\setminus\left\{0\right\},c\in A^{\times},\left(n_{i}\right)_{i\in I}\in\N^{\left(I\right)}$,其中$x=c\prod\limits_{i\in I}p_{i}^{n_{i}}$.
	\begin{enumerate}
		\item $x\in A\iff\forall i\in I\left(n_{i}>0\right)$.
		\item 我们称$\left(c,\left(n_{i}\right)_{i\in I}\right)$是$x$在$\left(p_{i}\right)_{i\in I}$下的不可约因子分解.在滥用术语的情况下,也称公式$x=c\prod\limits_{i\in I}p_{i}^{n_{i}}$为$x$的不可约因子分解.
		\item 设$y\in A$的不可约因子分解为$y=d\prod\limits_{i\in I}p_{i}^{m_{i}}$,则
		\begin{itemize}
			\item $x\mid y\iff\forall i\in I\left(n_{i}\leq m_{i}\right)$;
			\item $x$和$y$的一个最大公因子是$\prod\limits_{i\in I}p_{i}^{\sup\left(n_{i},m_{i}\right)}$,一个最小公倍数是$\prod\limits_{i\in I}p_{i}^{\sup\left(n_{i},m_{i}\right)}$.
		\end{itemize}
		\item 这个定理对于一种更特殊的PID成立,我们把它称为唯一分解整环(或UFD).以后会知道,多项式环和形式幂级数环都是UFD.
	\end{enumerate}
\end{remark}

